% =====================================================================
%  ПОРТФОЛІО ІНЖЕНЕРА-РОЗРОБНИКА ДРУКОВАНИХ ПЛАТ
%  Владислав Лисюк
%
%  Як зібрати:  xelatex portfolio.tex      (двічі — заради змісту/полів)
%  Або просто:  build.bat  /  make
%
%  Рушій ОБОВ'ЯЗКОВО XeLaTeX або LuaLaTeX (потрібен fontspec для кирилиці).
%  pdflatex НЕ підійде.
% =====================================================================

\documentclass[11pt,a4paper]{article}

% ---------- Поля сторінки (15 мм з усіх боків) ----------
\usepackage[a4paper,margin=15mm,bottom=18mm,footskip=8mm]{geometry}

% ---------- Шрифти та мова ----------
% Використано DejaVu Sans Condensed — є в системі, має повну кирилицю.
% Щоб змінити шрифт, просто впишіть іншу назву (напр. "Segoe UI" або "Calibri").
\usepackage{fontspec}
\setmainfont{DejaVu Sans Condensed}
\newfontfamily\headingfont{DejaVu Sans Condensed}

\usepackage{polyglossia}
\setdefaultlanguage{ukrainian}

% ---------- Графіка, кольори, допоміжні пакети ----------
\usepackage{graphicx}
\usepackage{xcolor}
\usepackage[most]{tcolorbox}
\usepackage{enumitem}          % компактні списки
\usepackage{microtype}
\usepackage{fancyhdr}
\usepackage{ragged2e}

% ---------- ПАЛІТРА (тут міняти кольори) ----------
\definecolor{accent}{HTML}{0F5C4A}   % акцент — темно-зелений
\definecolor{ink}{HTML}{14181D}      % основний текст
\definecolor{muted}{HTML}{5D6772}    % сірий для метаданих і підписів
\definecolor{chipbg}{HTML}{EEF3F1}   % заливка «чипів»
\definecolor{line}{HTML}{D7DDE3}     % тонкі лінії та рамки

\color{ink}

% ---------- Колонтитул: номер сторінки внизу праворуч ----------
\pagestyle{fancy}
\fancyhf{}
\renewcommand{\headrulewidth}{0pt}
\fancyfoot[R]{\footnotesize\color{muted}\thepage}

% Без відступів на початку абзацу, невеликий проміжок між абзацами
\setlength{\parindent}{0pt}
\setlength{\parskip}{4pt}

% Заборонити «висячі» рядки
\widowpenalty=10000
\clubpenalty=10000

% =====================================================================
%  ВЛАСНІ КОМАНДИ — далі ними описано весь документ
% =====================================================================

% ---- «Чип»: короткий тег у рамці зі світлою заливкою ----
\newtcbox{\chip}{on line, boxsep=0pt, left=5pt, right=5pt, top=3pt, bottom=3pt,
  colframe=line, colback=chipbg, boxrule=0.4pt, arc=3pt,
  fontupper=\footnotesize\color{accent!85!black}}

% ---- Шапка проєкту ----
% #1 — надзаголовок («ПРОЄКТ 01»), #2 — назва, #3 — рядок метаданих
\newcommand{\projecthead}[3]{%
  {\footnotesize\color{muted}\addfontfeature{LetterSpace=25}\MakeUppercase{#1}}\\[3pt]
  {\LARGE\bfseries\color{ink}#2}\\[4pt]
  {\footnotesize\color{muted}#3}\par
  \vspace{4pt}
  \hrulecolor{line}{0.6pt}
  \vspace{8pt}
}

% ---- Ряд «чипів»: \chips{\chip{...}\cs\chip{...}} ----
% (рядок вирівняно ліворуч, тож зайві чипи переносяться на новий рядок)
\newcommand{\chips}[1]{\par\noindent{\raggedright #1\par}\vspace{2pt}}

% ---- Проміжок між «чипами» ----
% (дозволяє перенести ряд чипів на новий рядок, якщо вони не вміщуються)
\newcommand{\cs}{\hskip3pt\penalty0\hskip0pt\relax}

% ---- Горизонтальна лінія на всю ширину ----
% #1 — колір, #2 — товщина
\newcommand{\hrulecolor}[2]{\par\noindent{\color{#1}\rule{\linewidth}{#2}}\par}

% ---- Блок «Моя роль»: вертикальна лінія акценту ліворуч ----
\newtcolorbox{rolebox}{enhanced, breakable=false, colback=white,
  boxrule=0pt, frame hidden, sharp corners,
  borderline west={2pt}{0pt}{accent},
  left=10pt, right=0pt, top=5pt, bottom=5pt, boxsep=0pt}

\newcommand{\role}[1]{%
  \par\vspace{6pt}
  \begin{rolebox}
    \small\justifying{\bfseries\color{accent}Моя роль.} #1
  \end{rolebox}
  \par\vspace{4pt}
}

% ---- ТЕЗИ під шапкою проєкту ----
% Короткі тези: під які вимоги плату спроєктовано / що саме адаптовано.
% Використання:
%   \notes{
%     \tz Перший рядок.
%     \tz Другий рядок.
%   }
\newcommand{\tz}{\item}
\newcommand{\notes}[1]{%
  \par\vspace{2pt}
  {\footnotesize\color{muted}%
   \begin{itemize}[leftmargin=12pt, itemsep=1pt, topsep=2pt, parsep=0pt,
     label=\textcolor{muted}{\rule{3pt}{3pt}}]
     #1
   \end{itemize}}%
  \par\vspace{2pt}
}

% ---- Марковані пункти ----
\newenvironment{points}
  {\begin{itemize}[leftmargin=14pt, itemsep=3pt, topsep=4pt, parsep=0pt,
     label=\textcolor{accent}{\small$\bullet$}]\small\justifying}
  {\end{itemize}}

% ---- ЗОБРАЖЕННЯ НА ВСЮ ШИРИНУ ----
% #1 — [необов'язково] максимальна висота (частка \textheight), типово 0.42
% #2 — файл, #3 — підпис
% Усе загорнуто в minipage: LaTeX ніколи не розриває minipage між сторінками,
% тож зображення разом із підписом завжди переходить на наступну сторінку цілком.
\newcommand{\bigimg}[3][0.42]{%
  \par\vspace{6pt}
  \noindent
  \begin{minipage}{\linewidth}
    \centering
    \setlength{\fboxrule}{0.4pt}\setlength{\fboxsep}{2pt}%
    \fcolorbox{line}{white}{\includegraphics[width=\dimexpr\linewidth-6pt\relax,
                                            height=#1\textheight,
                                            keepaspectratio]{#2}}\\[3pt]
    {\footnotesize\color{muted}#3}
  \end{minipage}
  \par\vspace{8pt}
}

% ---- ПАРА ЗОБРАЖЕНЬ У ОДИН РЯД ----
% #1 — [необов'язково] максимальна висота, типово 0.30
% #2/#3 — файл і підпис лівого; #4/#5 — файл і підпис правого
\newcommand{\pairimg}[5][0.30]{%
  \par\vspace{6pt}
  \noindent
  \begin{minipage}{\linewidth}
    \setlength{\fboxrule}{0.4pt}\setlength{\fboxsep}{2pt}%
    \begin{minipage}[t]{0.485\linewidth}
      \centering
      \fcolorbox{line}{white}{\includegraphics[width=\dimexpr\linewidth-6pt\relax,
                                              height=#1\textheight,
                                              keepaspectratio]{#2}}\\[3pt]
      {\footnotesize\color{muted}#3}
    \end{minipage}\hfill
    \begin{minipage}[t]{0.485\linewidth}
      \centering
      \fcolorbox{line}{white}{\includegraphics[width=\dimexpr\linewidth-6pt\relax,
                                              height=#1\textheight,
                                              keepaspectratio]{#4}}\\[3pt]
      {\footnotesize\color{muted}#5}
    \end{minipage}
  \end{minipage}
  \par\vspace{8pt}
}

\usepackage{calc}   % потрібен для \heightof у блоці «Моя роль»

\graphicspath{{images/}}

% =====================================================================
%  ДОКУМЕНТ
% =====================================================================
\begin{document}

% Документ починається одразу з першого проєкту — без шапки з контактами
% та без вступного абзацу.
% Проєкти впорядковані хронологічно: від першої роботи до останньої.

% ---------------------------------------------------------------
%  ПРОЄКТ 01 — термометр на ATtiny13 (курсовий проєкт)
% ---------------------------------------------------------------
\projecthead{Проєкт 01}%
  {Термометр на ATtiny13}%
  {Курсовий проєкт із дисципліни «Сучасні САПР РЕА» $\cdot$ 2 шари $\cdot$ 2023}

\chips{\chip{Altium Designer}\cs\chip{2 шари}\cs\chip{ATtiny13}\cs\chip{Перша робота}}

\notes{
  \tz Перша самостійна робота в Altium Designer.
  \tz Призначення - кімнатний термометр: датчик LM35, зсувні регістри 74HC595,
      три 7-сегментні індикатори.
}

\role{Бібліотека компонентів, схема електрична принципова,  трасування
 плати, комплект конструкторської документації.}

\bigimg[0.28]{temp_iso.jpg}{3D-вигляд в Altium Designer.}
\pairimg[0.30]{temp_3d_top.jpg}{Вигляд зверху.}%
              {temp_layout.jpg}{Розводка, два шари.}

% ---------------------------------------------------------------
%  ПРОЄКТ 02 — плата віддаленої лабораторії силової електроніки
% ---------------------------------------------------------------
\clearpage
\projecthead{Проєкт 02}%
  {Плата віддаленої лабораторії силової електроніки}%
  {Horizon Europe $\cdot$ NGI Search, грант № 101069364 $\cdot$ 2 шари $\cdot$ партія 5 шт $\cdot$ 2023--2024}

\chips{\chip{Altium Designer}\cs\chip{Понижувальний перетворювач}\cs\chip{Horizon Europe}}

\notes{
  \tz Перший виробничий проєкт в якому брав участь - плата пройшла повний шлях від схеми до
      виготовлення й монтажу.
}

\role{Бібліотека компонентів та оформлення схем в Altium Designer, компонування та
трасування плати, повне складання й пайка всіх п'яти плат.}


\bigimg[0.44]{buck_3d.jpg}{3D-вигляд в Altium Designer.}
\bigimg[0.40]{buck_layout.jpg}{Розводка.}
\bigimg[0.40]{ngi_assembled.jpg}{Зібрана плата.}
\pairimg[0.64]{ngi_batch.jpg}{Виготовлена партія.}%
              {ngi_stand.jpg}{У складі стенда разом із Raspberry Pi та Analog Discovery.}

\vspace{4pt}
{\footnotesize\color{muted}Виконано в рамках проєкту NGI Search, грантова угода № 101069364.}
\vspace{3pt}
\hrulecolor{line}{0.6pt}

% ---------------------------------------------------------------
%  ПРОЄКТ 03 — шилд для NUCLEO-H723
% ---------------------------------------------------------------
\clearpage
\projecthead{Проєкт 03}%
  {Плата розширення для STM32 NUCLEO-H723}%
  {Віддалена лабораторія мікропроцесорної техніки $\cdot$ 2 шари, висока щільність $\cdot$ Erasmus+ DIGITRANS $\cdot$ 2024--2025}

\chips{\chip{Altium Designer}\cs\chip{2 шари, висока щільність}\cs\chip{STM32H723}\cs%
\chip{C для МК}}

\notes{
  \tz Розведено лише на двох шарах відповідно до виробничого бюджету проєкту.
  \tz Посадкові місця й контур узгоджені з налагоджувальним стендом NUCLEO-H723.
  \tz П'ять комплектів зібрано вручну (THT + SMD) та інтегровано в навчальні стенди
      разом із Raspberry Pi і Analog Discovery.
}

\role{Бібліотека компонентів, робота зі схемами та розводка плати виконана під керівництвом провідного інженера.
Складання та пайка всіх п'яти комплектів, тестові прошивки.}



\pairimg[0.38]{shield_iso.jpg}{3D-вигляд в Altium Designer.}%
              {shield_layout.jpg}{Розводка, обидва шари.}
\pairimg[0.38]{shield_assembled.jpg}{Плата після SMD-монтажу.}%
              {shield_nucleo.jpg}{Встановлено на NUCLEO-H723.}
\bigimg[0.48]{shield_stand.jpg}{Стенд у робочому стані: Analog Discovery 2, шилд із Nucleo та Raspberry Pi на спільній панелі.}

% ---------------------------------------------------------------
%  ПРОЄКТ 04 — вимірювання температури термопарами
% ---------------------------------------------------------------
\clearpage
\projecthead{Проєкт 04}%
  {24-канальна система вимірювання температури термопарами}%
  {Промислова вимірювальна плата $\cdot$ 4 шари $\cdot$ STM32F4 $\cdot$ виготовлено в Німеччині $\cdot$ 2025--2026}

\chips{\chip{Altium Designer}\cs\chip{STM32}\cs\chip{Гальванічна розв'язка}\cs%
\chip{Термопари J/K}}

\notes{
  \tz Не розробка з нуля - адаптація наявного промислового проєкту під нові вимоги
      замовника.
}

\role{Перероблення схеми вимірювального тракту термопар, схема комутації каналів, вилучення невикористаних блоків проєкту. Зміни в бібліотеці компонентів, внесення змін в друковану плату.
Виконано під час стажування в Hochschule Bonn-Rhein-Sieg (DAAD Eastern Partnership).}



\bigimg[0.44]{tc_iso.jpg}{3D-вигляд в Altium Designer.}
\bigimg[0.40]{tc_top.jpg}{Вигляд зверху: функціональні зони.}
\bigimg[0.38]{tc_bottom.jpg}{Вид знизу.}

% ---------------------------------------------------------------
%  ПРОЄКТ 05 — драйвер електромагнітних клапанів
% ---------------------------------------------------------------
\clearpage
\projecthead{Проєкт 05}%
  {8-канальний драйвер електромагнітних клапанів}%
  {Силова електроніка $\cdot$ 4 шари $\cdot$ STM32 $\cdot$ виготовлено, партія 5 шт $\cdot$ 2026}

\chips{\chip{Altium Designer}\cs\chip{STM32}\cs\chip{Силова електроніка}}

\notes{
  \tz Частковий розрахунок схеми та розробка конструкції та плати під керівництвом провідного інженера.
  \tz Компонування підпорядковане вимозі розділення силової та керуючої частин:
      силові кола по периметру, мікроконтролер ізольовано в центрі.
}

\role{Частковий розрахунок номіналів схем, бібліотека компонентів, повна розводка друкованої плати. Робота під наглядом провідного інженера. Основа кваліфікаційної роботи бакалавра.}


\bigimg[0.46]{valve_3d_iso.jpg}{3D-вигляд в Altium Designer.}
\bigimg[0.38]{valve_3d_top.jpg}{Вигляд зверху: дзеркальна симетрія восьми силових каналів, керуюча частина в центрі.}
\bigimg[0.46]{valve_board.jpg}{Виготовлена плата (партія 5 шт) до монтажу компонентів.}

% ---------------------------------------------------------------
%  ПРОЄКТ 06 — плата, сумісна з Arduino Duemilanove
% ---------------------------------------------------------------
\clearpage
\projecthead{Проєкт 06}%
  {Плата, сумісна з Arduino Duemilanove}%
  {Навчальний проєкт $\cdot$ 2 шари $\cdot$ ATmega328 $\cdot$ THT + SMD}

\chips{\chip{Altium Designer}\cs\chip{2 шари}\cs\chip{ATmega328}\cs%
\chip{USB}\cs\chip{Живлення 5\,В}\cs\chip{Змішаний монтаж}}

\notes{
  \tz Навчальний проєкт: відтворення плати, сумісної з Arduino Duemilanove - 
      ATmega328 у DIP-корпусі, USB-міст.
  \tz Змішаний монтаж підібрано під ручну пайку: SMD на верхньому шарі,
      мікроконтролер і роз'єми - THT.
}

\role{Створення бібліотеки компонентів та розводка плати, навчальний проєкт.}


\pairimg[0.32]{ard_iso.jpg}{3D-вигляд в Altium Designer.}%
              {ard_top.jpg}{Вигляд зверху.}
\bigimg[0.32]{ard_layout.jpg}{Розводка, обидва шари.}

% ---------------------------------------------------------------
%  ПРОЄКТ 07 — драйвер затвора SiC MOSFET
% ---------------------------------------------------------------
\clearpage
\projecthead{Проєкт 07}%
  {Драйвер затвора SiC MOSFET для системи динамічного бездротового живлення}%
  {Модуль для інвертора DWPT $\cdot$ 2 шари $\cdot$ 2026}

\chips{\chip{Altium Designer}\cs\chip{SiC MOSFET}\cs\chip{Гальванічна розв'язка}\cs%
\chip{ANSYS}}

\notes{
  \tz Схема не власна: за основу взято референсний драйвер Wolfspeed CGD15SG00D2
      для SiC MOSFET 3-го покоління (C3M).
  \tz Конструкцію плати перероблено під вимоги інвертора гібридної системи
      динамічного бездротового передавання енергії - форму й посадкові розміри
      адаптовано під модульну заміну.
  \tz Драйвер виконано окремим модулем, що впаюється у плату інвертора: так вузол
      можна замінити, не чіпаючи основну плату.
  \tz Оптимізовано трасування вихідного каскаду для зменшення
      паразитних ємності та індуктивності при швидкому перемиканні.
}

\role{Адаптація й повне перекомпонування друкованої плати, бібліотека компонентів,
трасування вихідного каскаду, моделювання паразитних параметрів в ANSYS.}

\bigimg[0.44]{sic_iso.jpg}{3D-вигляд в Altium Designer.}
\pairimg[0.34]{sic_top.jpg}{Вид спереду.}%
              {sic_bottom.jpg}{Нижній бік: цифровий ізолятор і кола затвора.}
\bigimg[0.50]{sic_layout.jpg}{Розводка: контур плати підігнано під посадку в плату інвертора.}

\hrulecolor{line}{0.6pt}

\end{document}
